% Pista ana-26j-q1 · Anàlisi
% Procedència: PAU Catalunya 2026 · Convocatòria de juny · Sèrie 1 · Exercici 1
\documentclass[11pt,a4paper]{article}
\usepackage{pista}
\pistainfo{Catalunya 2026}{Convocatòria de juny}{Sèrie 1}

\begin{document}

\section*{\textcolor{blauPista}{Exercici 1}}

\noindent Considereu la funció definida a trossos següent:
\[
f(x) = \begin{cases} 5e^{2x} & x \le 0 \\ (x + m)^{2} + 1 & 0 < x < 2 \\ 1 & 2 \le x \end{cases},
\]
on $m$ és un paràmetre real.

\subsection*{1.1. \normalfont Determineu els valors de $m$ que fan que la funció $f(x)$ sigui contínua en tot el seu domini. Justifiqueu la resposta. \hfill\textnormal{[1 punt]}}

\pista{Cadascun dels tres trossos ja és continu dins del seu propi interval (són funcions elementals), així que només cal vigilar els dos punts de contacte, $x = 0$ i $x = 2$. A cadascun d'ells, iguala el límit per l'esquerra, el límit per la dreta i el valor de la funció en aquell punt (que hauria de coincidir amb un dels dos límits laterals, ja que és una funció a trossos amb desigualtats no estrictes). T'apareixeran dues condicions sobre $m$: resol-les per separat i queda't només amb els valors que en compleixen totes dues alhora.}

\subsection*{1.2. \normalfont Feu un esbós de la gràfica de $y = f(x)$ per al cas $m = -2$, i calculeu l'àrea delimitada per aquesta gràfica, l'eix $OX$ i les rectes $x = -1$ i $x = 3$. \hfill\textnormal{[1 punt]}}

\pista{Amb $m = -2$, el tros del mig és $(x - 2)^{2} + 1$, una paràbola amb el vèrtex desplaçat; identifica'l per fer l'esbós (o calcula'l amb la fórmula del vèrtex). Un cop vegis que la funció és sempre positiva a l'interval $[-1, 3]$, l'àrea que et demanen és, simplement, la integral definida de $f(x)$ entre $-1$ i $3$; com que $f$ ve definida a trossos, hauràs de descompondre aquesta integral en tres integrals més petites (una per cada tram on canvia l'expressió de $f$) i sumar-ne els resultats.}

\subsection*{1.3. \normalfont Per a $m = -2$, trobeu un punt on la recta tangent a $y = f(x)$ sigui paral·lela a $y = -2x$. Calculeu l'equació d'aquesta recta tangent. \hfill\textnormal{[0,5 punts]}}

\pista{Pensa en la forma de la gràfica que has fet a l'apartat anterior: el primer tram és creixent (exponencial) i el tercer és constant (pendent zero), de manera que l'única banda on la funció pot tenir pendent negatiu és al tram parabòlic del mig. Deriva aquest tram i iguala la derivada a $-2$ (el pendent de la recta donada). Un cop trobat el punt, aplica la fórmula punt-pendent per obtenir l'equació de la tangent.}

\end{document}
